\subsection{Bron--Kerbosch: Clique Máximo (grupo compatible)}
\label{alg:4-7}

\graybold{Preguntas guía:}

\begin{itemize}
\item ¿Necesito encontrar el clique máximo o todos los cliques maximales (no ampliables) de un grafo no dirigido?
\item ¿El número de vértices es relativamente pequeño (aproximadamente $n \le 50$)?
\item ¿La solución requiere explorar todas las posibles extensiones de un conjunto de vértices compatible?
\end{itemize}

\graybold{Utilidad:} Encontrar todos los cliques maximales o el clique máximo de un grafo no dirigido mediante backtracking con poda. Es la técnica estándar para resolver problemas de máxima compatibilidad, grupos completamente conectados y optimización sobre grafos pequeños.

\graybold{Complejidad:} \bigO{3^{n/3}} en el peor caso con el pivote de Tomita (el que maximiza $|P \cap N(u)|$), que es óptimo para enumerar todos los cliques maximales; aquí el pivote se elige arbitrariamente, así que esa cota no está garantizada, aunque en la práctica igual reduce mucho las llamadas recursivas. Memoria \bigO{n^2}: cada nivel de la recursión crea conjuntos nuevos de tamaño \bigO{n}.

\graybold{Fundamento teórico:} Durante la búsqueda se mantienen tres conjuntos:
\begin{itemize}
\item $R$: vértices que pertenecen al clique actual.
\item $P$: vértices que todavía pueden añadirse al clique.
\item $X$: vértices ya procesados, utilizados para evitar generar soluciones repetidas.
\end{itemize}
Cuando tanto $P$ como $X$ quedan vacíos, el conjunto $R$ es un clique maximal (ya no se puede ampliar); el clique máximo es el mayor de todos ellos. Para reducir la cantidad de ramas exploradas se selecciona un pivote $u$ y únicamente se consideran los vértices de $P \setminus N(u)$, donde $N(u)$ representa los vecinos del pivote. Esta optimización disminuye significativamente el árbol de búsqueda.~\cite{bron1973}

\graybold{Ejercicio aplicado}

Dado un grafo no dirigido, determinar el tamaño del mayor clique del grafo. Cada vértice representa una persona y existe una arista cuando dos personas son compatibles. Se desea encontrar el grupo más grande donde todos los integrantes sean compatibles entre sí.

\graybold{I/O:} Un solo caso: $n$ y $m$ en la primera línea y luego $m$ líneas con las aristas, con vértices numerados de 0 a $n-1$; se lee línea a línea con \texttt{input()}. El algoritmo resulta adecuado para grafos pequeños, donde una búsqueda exhaustiva con poda es suficiente para obtener la solución óptima.

\graybold{Código solución}

\begin{lstlisting}[language=Python]
def bronk(R, P, X):
    global best

    if not P and not X:
        best = max(best, len(R))
        return

    if P or X:
        u = next(iter(P | X))
    else:
        u = None

    if u is None:
        candidates = list(P)
    else:
        candidates = list(P - graph[u])

    for v in candidates:
        bronk(
            R | {v},
            P & graph[v],
            X & graph[v]
        )

        P.remove(v)
        X.add(v)

n, m = map(int, input().split())
graph = [set() for _ in range(n)]
for _ in range(m):
    u, v = map(int, input().split())
    graph[u].add(v)
    graph[v].add(u)

best = 0
bronk(set(), set(range(n)), set())
print(best)
\end{lstlisting}

\graybold{Observación:} El algoritmo de Bron--Kerbosch explora únicamente conjuntos que pueden convertirse en un clique válido. La versión con pivote es la utilizada habitualmente en programación competitiva debido a que reduce considerablemente el número de llamadas recursivas. Aunque su complejidad en el peor caso sigue siendo exponencial, resulta extremadamente eficiente para grafos pequeños, siendo la solución estándar para problemas de clique máximo en ICPC.
